% GENERATED FILE. Edit canonical lessons, then run npm run generate:books.
% canonical-source-hash: fnv1a64:9f7d63c4b3e2713b
% canonical-lessons: KA-C433-prayanika, KA-C433-platpharm, KA-C433-bogi, KA-C433-hali, KA-C433-talupu

\chapter{At the Railway Station}
\label{ch:ka-at-the-railway-station}
\hlchaptermodality{433}

\begin{chapteropening}
\textbf{By the end of this chapter:} \emph{I can say a passenger, a railway platform, a railway coach, a rail track and to reach.}

Complete the last lesson of chapter 433: I can say a passenger, a railway platform, a railway coach, a rail track and to reach.
\end{chapteropening}

\section[\texorpdfstring{\kn{ಪ್ರಯಾಣಿಕ} (prayāṇika)}{ಪ್ರಯಾಣಿಕ (prayāṇika)}]{\kn{ಪ್ರಯಾಣಿಕ} (prayāṇika) — a passenger; a traveller}

\label{lesson:KA-C433-prayanika}



\begin{quote}
Before the new one: say the Kannada for chalk, then the Kannada for a ruler.
\end{quote}

\subsection*{You'll want to know: \kn{ಪ್ರಯಾಣಿಕ}}
\textbf{\kn{ಪ್ರಯಾಣಿಕ}} — \emph{prayāṇika} — \textquotedblleft{}a passenger; a traveller\textquotedblright{}.

\begin{quote}
\kn{ರೈಲು} \kn{ನಿಲ್ದಾಣದಲ್ಲಿ} \kn{ತುಂಬಾ} \kn{ಪ್ರಯಾಣಿಕರು} \kn{ಇದ್ದರು}.
\end{quote}

\begin{quote}
\emph{railu nildāṇadalli tumbā prayāṇikaru iddaru.}
\end{quote}

\begin{quote}
There were a lot of passengers at the railway station.
\end{quote}

Plural: \kn{ಪ್ರಯಾಣಿಕರು}, \emph{prayāṇikaru}.

Built on \kn{ಪ್ರಯಾಣ}, a journey, with -\kn{ಇಕ}, one who belongs to it: the one who makes the journey.

\begin{cousinweb}
\textbf{\kn{ಪ್ರಯಾಣಿಕ}} (\emph{prayāṇika}) is built from Sanskrit \emph{prayāṇa}, \textquotedblleft{}setting out, a journey\textquotedblright{} (CDIAL 8732), with the ending -\emph{ika}, \textquotedblleft{}one who belongs to\textquotedblright{}. The same \emph{prayāṇa} gives \kn{ಪ್ರಯಾಣ}, a journey.
\end{cousinweb}

\begin{grammarlens}[title={What you've built}]
One more word for getting around.
\end{grammarlens}

\subsection*{Your turn}
\begin{itemize}
\raggedright
  \item \emph{Say it:} \emph{prayāṇika}
  \item \emph{Say it:} \emph{prayāṇika}, once more
  \item \emph{Say it:} \emph{railu nildāṇadalli tumbā prayāṇikaru iddaru.} — There were a lot of passengers at the railway station.
  \item \emph{Recall:} say the Kannada for chalk, then the Kannada for a ruler, then say \emph{prayāṇika} again
\end{itemize}

\begin{tcolorbox}[breakable,colback=teal!4,colframe=teal!35!black,title={Before you move on}]
What does \kn{ಪ್ರಯಾಣಿಕ} mean? (\textquotedblleft{}A passenger; a traveller\textquotedblright{}.) Say it once more.
\end{tcolorbox}
\section[\texorpdfstring{\kn{ಪ್ಲಾಟ್ಫಾರ್ಮ್} (plāṭphārm)}{ಪ್ಲಾಟ್ಫಾರ್ಮ್ (plāṭphārm)}]{\kn{ಪ್ಲಾಟ್ಫಾರ್ಮ್} (plāṭphārm) — a railway platform}

\label{lesson:KA-C433-platpharm}



\begin{quote}
Before the new one: say the Kannada for a ruler, then the Kannada for a passenger.
\end{quote}

\subsection*{You'll want to know: \kn{ಪ್ಲಾಟ್ಫಾರ್ಮ್}}
\textbf{\kn{ಪ್ಲಾಟ್ಫಾರ್ಮ್}} — \emph{plāṭphārm} — \textquotedblleft{}a railway platform\textquotedblright{}.

\begin{quote}
\kn{ನಮ್ಮ} \kn{ರೈಲು} \kn{ಮೂರನೇ} \kn{ಪ್ಲಾಟ್ಫಾರ್ಮ್ಗೆ} \kn{ಬರುತ್ತದೆ}.
\end{quote}

\begin{quote}
\emph{namma railu mūranē plāṭphārmge baruttade.}
\end{quote}

\begin{quote}
Our train comes in at the third platform.
\end{quote}

\kn{ವೇದಿಕೆ} is a stage or any raised platform; at a station people say \kn{ಪ್ಲಾಟ್ಫಾರ್ಮ್}.

\begin{cousinweb}
\textbf{\kn{ಪ್ಲಾಟ್ಫಾರ್ಮ್}} (\emph{plāṭphārm}) is a loan from English \emph{platform}, which came in with the railways. Its last consonant keeps the virama, as recent English loans do.
\end{cousinweb}

\begin{grammarlens}[title={What you've built}]
One more word for getting around.
\end{grammarlens}

\subsection*{Your turn}
\begin{itemize}
\raggedright
  \item \emph{Say it:} \emph{plāṭphārm}
  \item \emph{Say it:} \emph{plāṭphārm}, once more
  \item \emph{Say it:} \emph{namma railu mūranē plāṭphārmge baruttade.} — Our train comes in at the third platform.
  \item \emph{Recall:} say the Kannada for a ruler, then the Kannada for a passenger, then say \emph{plāṭphārm} again
\end{itemize}

\begin{tcolorbox}[breakable,colback=teal!4,colframe=teal!35!black,title={Before you move on}]
What does \kn{ಪ್ಲಾಟ್ಫಾರ್ಮ್} mean? (\textquotedblleft{}A railway platform\textquotedblright{}.) Say it once more.
\end{tcolorbox}
\section[\texorpdfstring{\kn{ಬೋಗಿ} (bōgi)}{ಬೋಗಿ (bōgi)}]{\kn{ಬೋಗಿ} (bōgi) — a railway coach; a carriage}

\label{lesson:KA-C433-bogi}



\begin{quote}
Before the new one: say the Kannada for a passenger, then the Kannada for a railway platform.
\end{quote}

\subsection*{You'll want to know: \kn{ಬೋಗಿ}}
\textbf{\kn{ಬೋಗಿ}} — \emph{bōgi} — \textquotedblleft{}a railway coach; a carriage\textquotedblright{}.

\begin{quote}
\kn{ನನ್ನ} \kn{ಅಣ್ಣ} \kn{ಮುಂದಿನ} \kn{ಬೋಗಿಯಲ್ಲಿ} \kn{ಕುಳಿತಿದ್ದಾನೆ}.
\end{quote}

\begin{quote}
\emph{nanna aṇṇa mundina bōgiyalli kuḷitiddāne.}
\end{quote}

\begin{quote}
My elder brother is sitting in the next coach.
\end{quote}

Plural: \kn{ಬೋಗಿಗಳು}, \emph{bōgigaḷu}.

A train is a line of \kn{ಬೋಗಿಗಳು} pulled by an engine.

\begin{cousinweb}
\textbf{\kn{ಬೋಗಿ}} (\emph{bōgi}) is a loan from English \emph{bogie}, the wheeled frame under a railway carriage. In Indian use the word came to name the whole coach.
\end{cousinweb}

\begin{grammarlens}[title={What you've built}]
One more word for getting around.
\end{grammarlens}

\subsection*{Your turn}
\begin{itemize}
\raggedright
  \item \emph{Say it:} \emph{bōgi}
  \item \emph{Say it:} \emph{bōgi}, once more
  \item \emph{Say it:} \emph{nanna aṇṇa mundina bōgiyalli kuḷitiddāne.} — My elder brother is sitting in the next coach.
  \item \emph{Recall:} say the Kannada for a passenger, then the Kannada for a railway platform, then say \emph{bōgi} again
\end{itemize}

\begin{tcolorbox}[breakable,colback=teal!4,colframe=teal!35!black,title={Before you move on}]
What does \kn{ಬೋಗಿ} mean? (\textquotedblleft{}A railway coach; a carriage\textquotedblright{}.) Say it once more.
\end{tcolorbox}
\section[\texorpdfstring{\kn{ಹಳಿ} (haḷi)}{ಹಳಿ (haḷi)}]{\kn{ಹಳಿ} (haḷi) — a rail track; a rail}

\label{lesson:KA-C433-hali}



\begin{quote}
Before the new one: say the Kannada for a railway platform, then the Kannada for a railway coach.
\end{quote}

\subsection*{You'll want to know: \kn{ಹಳಿ}}
\textbf{\kn{ಹಳಿ}} — \emph{haḷi} — \textquotedblleft{}a rail track; a rail\textquotedblright{}.

\begin{quote}
\kn{ರೈಲು} \kn{ಹಳಿಯ} \kn{ಮೇಲೆ} \kn{ನಡೆಯಬಾರದು}.
\end{quote}

\begin{quote}
\emph{railu haḷiya mēle naḍeyabāradu.}
\end{quote}

\begin{quote}
Nobody should walk on the rail track.
\end{quote}

Plural: \kn{ಹಳಿಗಳು}, \emph{haḷigaḷu}.

\kn{ಹಳಿ} \kn{ತಪ್ಪು}, to leave the rails, is said of a train and, loosely, of a plan that goes wrong.

\begin{cousinweb}
\textbf{\kn{ಹಳಿ}} (\emph{haḷi}) looks like a native Kannada word, but its origin is not settled: DEDR shows no clear source for \emph{haḷi} in the sense of a rail or a track.
\end{cousinweb}

\begin{grammarlens}[title={What you've built}]
One more word for getting around.
\end{grammarlens}

\subsection*{Your turn}
\begin{itemize}
\raggedright
  \item \emph{Say it:} \emph{haḷi}
  \item \emph{Say it:} \emph{haḷi}, once more
  \item \emph{Say it:} \emph{railu haḷiya mēle naḍeyabāradu.} — Nobody should walk on the rail track.
  \item \emph{Recall:} say the Kannada for a railway platform, then the Kannada for a railway coach, then say \emph{haḷi} again
\end{itemize}

\begin{tcolorbox}[breakable,colback=teal!4,colframe=teal!35!black,title={Before you move on}]
What does \kn{ಹಳಿ} mean? (\textquotedblleft{}A rail track; a rail\textquotedblright{}.) Say it once more.
\end{tcolorbox}
\section[\texorpdfstring{\kn{ತಲುಪು} (talupu)}{ತಲುಪು (talupu)}]{\kn{ತಲುಪು} (talupu) — to reach; to arrive at}

\label{lesson:KA-C433-talupu}



\begin{quote}
Before the new one: say the Kannada for a railway coach, then the Kannada for a rail track.
\end{quote}

\subsection*{You'll want to know: \kn{ತಲುಪು}}
\textbf{\kn{ತಲುಪು}} — \emph{talupu} — \textquotedblleft{}to reach; to arrive at\textquotedblright{}.

\begin{quote}
\kn{ರೈಲು} \kn{ರಾತ್ರಿ} \kn{ಹತ್ತು} \kn{ಗಂಟೆಗೆ} \kn{ಊರು} \kn{ತಲುಪಿತು}.
\end{quote}

\begin{quote}
\emph{railu rātri hattu gaṇṭege ūru talupitu.}
\end{quote}

\begin{quote}
The train reached the town at ten at night.
\end{quote}

Past (I …): \kn{ತಲುಪಿದೆ}, \emph{talupide}.

The place may take -\kn{ಅನ್ನು} or no ending: \kn{ಊರನ್ನು} \kn{ತಲುಪು}, \kn{ಊರು} \kn{ತಲುಪು}. \kn{ತಲುಪಿಸು}, to deliver, makes something else reach.

\begin{cousinweb}
\textbf{\kn{ತಲುಪು}} (\emph{talupu}) is native Dravidian: DEDR 3102 lists Kannada \emph{talapu}, \emph{talupu}, \textquotedblleft{}to reach, to arrive\textquotedblright{}, beside \emph{talapisu}, \textquotedblleft{}to cause to reach\textquotedblright{}, the source of \kn{ತಲುಪಿಸು}.
\end{cousinweb}

\begin{grammarlens}[title={What you've built}]
One more word for getting around.
\end{grammarlens}

\subsection*{Your turn}
\begin{itemize}
\raggedright
  \item \emph{Say it:} \emph{talupu}
  \item \emph{Say it:} \emph{talupu}, once more
  \item \emph{Say it:} \emph{railu rātri hattu gaṇṭege ūru talupitu.} — The train reached the town at ten at night.
  \item \emph{Recall:} say the Kannada for a railway coach, then the Kannada for a rail track, then say \emph{talupu} again
\end{itemize}

\begin{tcolorbox}[breakable,colback=teal!4,colframe=teal!35!black,title={Before you move on}]
What does \kn{ತಲುಪು} mean? (\textquotedblleft{}To reach; to arrive at\textquotedblright{}.) Say it once more.
\end{tcolorbox}
